\chapter{Обоснование архитектуры аппаратно-программного комплекса контроля процесса печати}

Архитектура системы контроля 3D-печати связывает получение изображений, обнаружение дефектов и передачу результата в контур управления принтером. Физические причины и визуальные признаки отклонений определяют способ наблюдения и обработку данных. Выбор аппаратных компонентов и программных средств зависит от условий съемки, доступности вычислительных ресурсов и предусмотренной реакции на дефект.

При сравнении решений рассматриваются доступность кадров, связь результата с печатным заданием и возможность выполнить управляющую команду. Надежность обмена оценивается с учетом недоступности сервиса анализа и отсутствия ответа, а производительность~\textemdash{} по частоте обработки кадров и задержке реакции. Эти критерии задают основания для выбора архитектуры, оборудования и модели распознавания.

\section{Требования к системе обнаружения дефектов в контуре управления 3D-принтером и анализ существующих решений}

Обнаружение дефекта типа «спагетти» по изображениям веб-камеры исследовалось как задача классификации с помощью сверточной нейронной сети. Тестовая выборка включала~100 изображений, поровну представлявших печать с дефектом и без него. Общая точность классификации составила~96~\%. Среди изображений с дефектом правильно распознано~94~\%, среди кадров без дефекта~\textemdash{} 98~\%. При обнаружении нарушения авторы предусматривают уведомление оператора или остановку печати~\cite{kim2020image}. Эти результаты характеризуют распознавание отдельных изображений и служат основанием для выбора бинарной классификации как исходной задачи анализа.

Для включения такого анализа в контур печати результат должен быть связан с исходным кадром и текущим заданием. Решение об отмене должно приниматься после подтверждения дефекта по последовательности кадров, чтобы единичная ошибка классификации не прерывала исправную печать. После подтверждения критического дефекта локальная логика должна формировать запрос на отмену задания штатной командой управления.

Для наблюдения за 3D-печатью и обнаружения дефектов уже существуют готовые программные решения. OctoPrint~\textemdash{} программная платформа с веб-интерфейсом для удаленного наблюдения за состоянием 3D-принтера, просмотра изображений камеры и управления печатью. Ее возможности расширяются подключаемыми модулями. Obico представляет собой систему автоматического обнаружения дефектов по изображениям камеры, которая может использоваться совместно с OctoPrint и поддерживает уведомление оператора и остановку печати. Detector~2~\textemdash{} подключаемый модуль OctoPrint для анализа изображений и отправки предупреждений при обнаружении отклонений~\cite{gabstur2024detection}.

При испытании системы на основе OctoPrint с модулями Obico и Detector~2 в отдельных случаях Obico обнаруживал нарушение только после завершения печати. Среди возможных причин ограничений работы системы авторы рассматривают фокусировку камеры, освещение и положение изделия~\cite{gabstur2024detection}. Различия освещения, оборудования и изображений изделий также влияют на качество моделей мониторинга~\cite{zhang2024domain}. Следовательно, результаты готового средства в одной конфигурации не определяют своевременность обнаружения дефекта на другом принтере.

Выбор собственного модуля анализа связан с адаптацией модели к экспериментальному принтеру и раздельной настройкой распознавания и подтверждения дефекта. Для адаптации предусматриваются обучение на изображениях рабочей камеры, разбор ошибочных срабатываний и дообучение при изменении условий съемки. Обновление модели при этом не требует изменения способа передачи команды отмены. Сохраненные кадры и ответы дают возможность установить, как оценка модели и правило подтверждения привели к управляющему действию.

Условия съемки входят в границы применимости модели. В исследовании визуального контроля поверхности параметры предварительной обработки подбирались при фиксированном освещении~\cite{измайлов2020анализ}. На экспериментальном принтере ракурс, геометрия рабочей зоны, форма и цвет изделий могут менять видимость дефекта и сходство его признаков с фоном~\cite{gabstur2024detection,zhang2024domain}. Поэтому проверочные данные должны включать разные условия наблюдения, а также кадры исправной печати с элементами конструкции, похожими на признаки дефекта.

Анализ на отдельном сервере в локальной сети освобождает хост принтера от выполнения нейросетевой модели, но делает мониторинг зависимым от сетевого обмена. Задержки передачи и недоступность сервера учитываются как ограничения выбранной архитектуры. Исходный кадр должен сохраняться до отправки, а ошибка обмена~\textemdash{} фиксироваться вместе со сведениями о задании. Отсутствие достоверного ответа считается неопределенным состоянием и не вызывает управляющего действия. Ответ от завершенного или предыдущего задания также не служит основанием для отмены текущей печати. Для разбора срабатываний необходимо сохранять время получения кадра и ответа, состояние задания, версию модели и действовавшие параметры анализа.

\section{Выбор аппаратной платформы управления, камеры и вычислительного узла обработки данных}

Аппаратная конфигурация строится вокруг трех узлов: хоста принтера, микроконтроллерной платы и сервера анализа. Камера подключается к хосту, который передает изображения на сервер. Такое распределение определяет требования к подключениям и вычислительным ресурсам каждого узла.

Для управления выбран Klipper~\textemdash{} программное обеспечение 3D-принтера, распределяющее задачи между управляющим компьютером и микроконтроллерной платой. Компьютер, далее называемый хостом, обрабатывает команды и рассчитывает движение, а микроконтроллер управляет исполнительными узлами в заданные моменты времени. Сценарии работы принтера задаются макросами. Сетевой обмен организуется через Moonraker~\textemdash{} сервис с программным интерфейсом для передачи команд Klipper и получения состояния принтера. Этот интерфейс дает модулю контроля доступ к состоянию текущего задания и штатной команде отмены, сохраняя единый способ управления при замене модели обнаружения.

В качестве хоста выбран одноплатный компьютер Raspberry Pi~3B. На нем предусматриваются работа Klipper и Moonraker, периодическое получение кадров и обмен с сервером. Для этой роли используются подключение камеры, сетевой интерфейс и связь с микроконтроллерной платой. Вынос нейросетевой модели на сервер исключает ее вычислительные требования из требований к хосту. Проверка достаточности ресурсов Raspberry Pi~3B для совместного выполнения этих задач предусмотрена при испытании конфигурации.

Для визуального мониторинга выбрана камера Sonix USB~2.0, подключаемая к Raspberry Pi~3B через универсальную последовательную шину (Universal Serial Bus, USB). Изображения поступают непосредственно на хост, откуда передаются на анализ. Камера должна закрепляться в постоянном положении относительно рабочей зоны, чтобы сохранять ракурс наблюдения на протяжении задания. При компоновке учитываются обзор стола и растущего изделия, перекрытие сцены подвижными узлами и влияние вибраций на резкость. Экспозиция и фокусировка выбираются по различимости признаков дефекта на разных этапах печати. Смещение камеры или изменение поля зрения меняет входные данные модели, поэтому требует повторной оценки качества распознавания.

На сервер в локальной сети возлагаются предварительная обработка изображений и запуск нейросетевой модели. Он возвращает результат на Raspberry Pi~3B, где локальный модуль проверяет принадлежность ответа активному заданию и применяет правило подтверждения дефекта. Управляющее действие передается через Moonraker. Сервер не управляет двигателями и нагревателями непосредственно. Такое распределение закрепляет вычислительные задачи за сервером, а принятие решения и взаимодействие с исполнительной частью~\textemdash{} за хостом.

Интервал получения кадров определяет время от появления видимого дефекта до первого изображения с его признаками. Задержка реакции дополнительно включает передачу данных, выполнение модели и подтверждение по последовательности кадров. Уменьшение интервала сокращает ожидание очередного снимка, но увеличивает число запросов к серверу. Поэтому выбор частоты связан с длительностью обработки и скоростью обмена. В испытаниях интервал, время ответа и число кадров для подтверждения рассматриваются совместно, чтобы объяснить итоговую задержку отмены.

Для управления исполнительной частью выбрана MKS Robin Nano~V1.2~\textemdash{} микроконтроллерная плата 3D-принтера, к которой подключаются драйверы шаговых двигателей, нагреватели и датчики. Микроконтроллер формирует шаговые импульсы с заданными временными интервалами независимо от получения кадров и сетевых задач хоста. Связь с Raspberry Pi~3B предусматривается через универсальный асинхронный приемопередатчик (Universal Asynchronous Receiver Transmitter, UART). По этой линии результаты расчета движения в Klipper передаются для выполнения подключенными узлами.

\section{Выбор инструментов разметки данных и подхода к обучению модели распознавания}

Для формирования набора данных планируется извлекать кадры из ускоренных видеозаписей печати, составленных из периодических снимков (таймлапсов). Источником служат записи с принтеров компании Stereotech, в том числе модели Hybrid~530~5.2.1. Эти изображения показывают изменение вида изделия в ходе задания. Открытый размеченный набор изображений предполагается использовать для расширения разнообразия сцен. Сведения о происхождении кадров сохраняются, поскольку соответствие обучающих изображений условиям рабочей камеры определяет применимость модели к экспериментальному принтеру.

Исходная схема разметки предусматривает ограничивающие прямоугольники для четырех классов: растрескивание, отлипание модели от стола, дефект типа «спагетти» и смещение слоев. Дополнительный класс «деталь» обозначает область изделия. Объектная детекция определяет положение обнаруженной области относительно печатаемой модели~\cite{брехт2022применение}. Такая разметка сохраняет возможность отдельных экспериментов с локализацией, хотя для первого обучения выбрана бинарная классификация наличия «спагетти». Аналогичная постановка с разделением изображений исправной и дефектной печати использована в работе~\cite{kim2020image}.

Для разметки выбран инструмент аннотирования изображений для задач компьютерного зрения (Computer Vision Annotation Tool, CVAT) в общедоступной редакции Community. CVAT поддерживает разметку ограничивающими прямоугольниками, единую схему классов и экспорт аннотаций для обучения. Предусматривается его локальное развертывание. Качество разметки оценивается по правильности класса, границам области и отсутствию повторных рамок одного дефекта. Перед включением аннотаций в набор проверяются соответствие идентификаторов классов принятой схеме и корректность координат.

Для каждой версии набора фиксируются источник изображений, состав аннотаций и контрольные суммы файлов. Это связывает исходные кадры с разметкой и сохраняет возможность повторного обучения на том же составе данных. Координаты рамок остаются в исходных аннотациях для экспериментов с детекторами.

Для обучения, валидации и тестирования изображения должны разделяться по сессиям печати. Все кадры одного задания должны относиться к одной части набора: соседние изображения часто имеют почти одинаковую композицию, и их распределение по разным частям способно завысить оценку качества. Проверочные данные должны охватывать разные изделия, материалы и условия освещения, встречающиеся на выбранном принтере.

Для определения наличия «спагетти» выбрана ResNet18~\textemdash{} сверточная нейронная сеть с остаточными связями (Residual Network, ResNet) в варианте с~18 слоями. Остаточные связи передают вход блока к его выходу в обход части слоев и облегчают обучение. Этот вариант требует меньше вычислений, чем более глубокие сети того же семейства~\cite{he2016deep}, что служит основанием для его выбора как базового классификатора. Предусматриваются веса, предварительно обученные на ImageNet~\textemdash{} наборе размеченных изображений для распознавания объектов. Исходный выходной слой планируется заменить одним нейроном и дообучить все слои на подготовленных кадрах. На вход планируется подавать полное изображение зоны печати, приведенное к единому размеру с сохранением пропорций.

Для отдельного выделения области детали предусмотрен RF-DETR Small~\textemdash{} детектор компании Roboflow на основе трансформера для обнаружения объектов (Detection Transformer, DETR)~\cite{robinson2026rfdetr}. Он возвращает координаты рамки изделия. Качество локализации оценивается независимо от распознавания «спагетти», а полученная рамка не используется для обрезки входа классификатора. Полный кадр сохраняет окружающее пространство, в котором также могут находиться признаки дефекта.

Метка кадра для обучения ResNet18 формируется по наличию целевого класса в исходных аннотациях независимо от числа размеченных областей. Набор планируется разделить на обучающую, валидационную и тестовую части в соотношении~60/20/20 с учетом известных сессий и семейств изображений. Обучающая часть используется для изменения весов, валидационная~\textemdash{} для выбора лучшей версии модели и порога срабатывания, тестовая~\textemdash{} для итоговой оценки. При преобладании отрицательных кадров состав каждой обучающей эпохи предполагается балансировать в соотношении~50/50. Контрольные выборки должны сохранять полный состав без такого выравнивания.

План проверки классификатора включает оценку общей доли верных ответов, точности положительных ответов, полноты и гармонического среднего точности и полноты (меры F1). Пропуски дефекта и ошибочные ответы на кадрах исправной печати учитываются отдельно, поскольку они приводят к разным последствиям для управления. Числовая оценка преобразуется в признак наличия дефекта сравнением с порогом, выбираемым на валидационных данных. Итоговая оценка должна выполняться на тестовой выборке без повторной настройки.

Проверка интеграции планируется на новых печатных сессиях с рабочей камеры. Здесь единицей оценки становится задание: рассматриваются пропуски дефекта в ходе печати, ошибочные отмены исправных заданий и задержка от появления видимого дефекта до выполнения отмены. Сохраненная последовательность ответов связывает ошибки классификации с результатом подтверждения. Проверка должна охватывать получение кадров, серверный анализ и выполнение команды как единый процесс. Такой подход отделяет качество модели на изображениях от качества реакции принтера в ходе задания.
